\section{Designs Under Test}
\label{sec:designs}

Every design serves \texttt{GET /result?roll=$n$} with the same JSON body, uses
the same synthetic data, and gets the same budget of two CPU cores for the read
path (Table~\ref{tab:designs}). Before the announced release time every design
answers ``not released'' without touching result data.

\begin{table}[t]
\caption{Designs and their share of the 2-core budget}
\label{tab:designs}
\centering
\footnotesize
\setlength{\tabcolsep}{3pt}
\begin{tabular}{@{}lp{4.6cm}l@{}}
\toprule
ID & Read path & CPUs \\
\midrule
A & App queries PostgreSQL (4-table join), builds the page & 1 + 1 \\
M & Two independent copies of A (``mirror links''); students pick one and may switch after an error & 2$\times$(0.5+0.5) \\
B0 & Redis cache-aside in front of A & 0.75+0.75+0.5 \\
B1 & B0 + request coalescing per key & same \\
B2 & B0 + XFetch~\cite{vattani2015xfetch} + 10\% TTL jitter & same \\
B3 & B0 with the cache filled before release & same \\
C1 & All results pre-computed into Redis; no database on the read path & 1 + 1 \\
C2 & One pre-computed JSON file per student, served by Nginx & 2 \\
D-A & Waiting room in front of A & 0.25 + 0.875 + 0.875 \\
D-C1 & Waiting room in front of C1 & 0.25 + 0.875 + 0.875 \\
\bottomrule
\end{tabular}
\end{table}

\textbf{Server model.} All application designs run the same Go server with a
pool of 256 workers and a queue of 4096 waiting requests, beyond which it
answers 503, as PHP-FPM or Apache would. Without such a bound, our
database-backed portal was killed for exceeding its memory limit during
overload. Building one page from the database costs a configurable amount of
CPU that stands for framework, session and rendering overhead; we use
\SI{5}{ms}. A Go server alone spends microseconds per page, which would make
design~A far cheaper than real portal stacks. Pre-computed designs store the
finished page and never pay this cost; cache-aside pays it on every miss.

\textbf{Waiting room.} Our waiting room admits students at a fixed rate (80\%
of the backend's calibrated capacity) and tells the others their position.
Two queue policies are compared: first-come-first-served (FIFO), and a
lottery that shuffles everyone who arrived before release, as some products
offer~\cite{cloudflare2026wr}. Two ticket implementations are compared: a
\emph{naive} ticket held only in the page (lost when the student presses
refresh) and a \emph{token} kept in a signed cookie (kept on refresh).
Tickets whose holder has not polled for 30\,s are skipped.
